% from main.tex
%%%%%%%%%%  scan idx 273  |  folio 260  %%%%%%%%%%
foncteur \(M\mapsto H^p(E\otimes^{\mathbf L}M)\) de \(\Mod(X)\) dans \(\Mod(X)\) est exact à droite.
\item[a')] La flèche canonique
\[
 H^p(E)\otimes M\longrightarrow H^p(E\otimes^{\mathbf L}M)
\]
est un isomorphisme pour tout \(M\in\Ob\Mod(X)\).
\item[b)] Le foncteur \(M\mapsto H^{p+1}(E\otimes^{\mathbf L}M)\) de \(\Mod(X)\) dans \(\Mod(X)\) est exact à gauche.
\item[b')] La flèche canonique définie ci-dessous
\begin{equation}
 H^{p+1}(E\otimes^{\mathbf L}M)\longrightarrow\underline{\Hom}(H^{-p-1}(E^{\vee}),M)\tag{$*$}
\end{equation}
est un isomorphisme pour tout \(M\in\Ob\Mod(X)\).
\end{enumerate}
Alors on a :
\[
 \mathrm a)\Longleftrightarrow\mathrm a')\Longleftrightarrow\mathrm b)\Longleftarrow\mathrm b').
\]
En outre, si \(E\) est pseudo-cohérent, les conditions a), a'), b), b') sont équivalentes.

\emph{Définition de} \((*)\) : La flèche canonique (I 7.5) \(E^{\vee}\otimes^{\mathbf L}E\longrightarrow\mathcal O_X\), où \(E^{\vee}=\mathrm R\underline{\Hom}(E,\mathcal O_X)\), donne, par tensorisation avec \(M\in\Ob\Mod(X)\), une flèche \(E^{\vee}\otimes^{\mathbf L}E\otimes^{\mathbf L}M\longrightarrow M\), d'où par la formule d'adjonction chère à Cartan (I 7.4), une flèche canonique
\begin{equation}
 E\otimes^{\mathbf L}M\longrightarrow\mathrm R\underline{\Hom}(E^{\vee},M).\tag{$**$}
\end{equation}
Composant la flèche induite par \((**)\) sur \(H^{p+1}\) avec la flèche canonique
\begin{equation}
 \underline{\Ext}^{p+1}(E^{\vee},M)\longrightarrow\underline{\Hom}(H^{-p-1}(E^{\vee}),M),\tag{$***$}
\end{equation}
on obtient la flèche \((*)\).

\emph{Preuve.} Les implications
\[
 \mathrm a')\Longrightarrow\mathrm a)\Longleftrightarrow\mathrm b)\Longleftarrow\mathrm b')
\]
sont triviales. Prouvons \(\mathrm a)\Longrightarrow\mathrm a')\). Soit \(M\) un \(\mathcal O_X\)-Module. Il existe une suite exacte \(M^1\longrightarrow M^0\longrightarrow M\longrightarrow0\), où \(M^0\) et \(M^1\) sont sommes directes de faisceaux du
